% answers_short.tex —— 2009 年·普通高等学校招生考试·山东卷·文科 · 简短答案速查
% 由 tools/make_short_answers.py 自 answers.tex 的 \gksolhead/\gkanswer 提取，
% 并与 2009/2009_数学真题_山东文_答案要点.md 逐题比对后生成，勿手改（改源 MD 后重跑）。

\begin{gknotice}
\textbf{使用说明}：本节只给最终答案便于快速核对，完整推导见解析版。答案为编写者独立推导，\textbf{非官方标准答案}。
\end{gknotice}

\gkpart{选择题}
\begin{center}\begin{tabular}{ccccccccccccc}
\toprule
题号 & 1 & 2 & 3 & 4 & 5 & 6 & 7 & 8 & 9 & 10 & 11 & 12 \\
\midrule
答案 & D & C & B & C & B & A & B & C & B & B & A & D \\
\bottomrule
\end{tabular}\end{center}

\gkpart{填空题}
\begin{itemize}[leftmargin=2.6em]
  \item[\textbf{13.}] $13$。
  \item[\textbf{14.}] $(1,+\infty)$。
  \item[\textbf{15.}] $30$。
  \item[\textbf{16.}] $2300$。
\end{itemize}

\gkpart{解答题}
\begin{itemize}[leftmargin=2.6em]
  \item[\textbf{17.}] U（原页条件与所给函数不相容，无法确定 $\varphi$）。
  \item[\textbf{18.}] （1）（2）证明见解析。
  \item[\textbf{19.}] （1）$z=400$；（2）$\frac7{10}$；（3）$\frac34$。
  \item[\textbf{20.}] （1）$r=-1$；（2）$T_n=\frac32-\frac{n+3}{2^{n+1}}$。
  \item[\textbf{21.}] （1）$b^2>a$；（2）$0<a<1$ 时 $b\geqslant-\frac{a+1}{2}$；$a\geqslant1$ 时 $b\geqslant-\sqrt a$。
  \item[\textbf{22.}] （1）$mx^2+y^2=1$；$m=0$ 时为两条平行直线，$m>0$ 时为椭圆，$m<0$ 时为双曲线；（2）$x^2+y^2=\frac45$；（3）$R=\sqrt2$ 时 $|A_1B_1|$ 最大，最大值为 $1$。
\end{itemize}
